\medskip

\textbf{Partie A}

\medskip

On considère deux cubes ABCDEFGH et BKLCFJMG positionnés comme sur la figure suivante:

\begin{center}
\psset{unit=1cm,arrowsize=2pt 3}
\begin{pspicture}(10.2,7.8)
\pspolygon(0.2,2.5)(5.8,0.3)(5.8,4.1)(0.2,6.3)
\psline(3,1.4)(3,5.2)(6.8,6.2)
\psline(5.8,0.3)(9.6,1.3)(9.6,5.1)(5.8,4.1)
\psline(6.8,6.2)(4,7.3)(0.2,6.3)
\psline(9.6,5.1)(6.8,6.2)
\psline[linestyle=dashed](4,7.3)(4,3.5)(9.6,1.3)
\pspolygon[linestyle=dashed](3,1.4)(6.8,2.4)(6.8,6.2)
\psline[linestyle=dashed](0.2,2.5)(4,3.5)

\psline[linewidth=1.25pt]{->}(0.2,2.5)(3,1.4)
\psline[linewidth=1.25pt,linestyle=dashed]{->}(0.2,2.5)(4,3.5)
\psline[linewidth=1.25pt]{->}(0.2,2.5)(0.2,6.3)
\psline(3,1.4)(1.6,5.75)(6.8,6.2)
\uput[dl](0.2,2.5){A} \uput[dl](3,1.4){B} \uput[ur](6.8,2.4){C} \uput[ul](4,3.5){D}
\uput[l](0.2,6.3){E} \uput[dl](3,5.2){F} \uput[ur](6.8,6.2){G} \uput[u](4,7.3){H}
\uput[u](5.8,4.1){J} \uput[d](5.8,0.3){K} \uput[r](9.6,1.3){L} \uput[ur](9.6,5.1){M}
\uput[u](1.6,5.75){I}
\end{pspicture}
\end{center}

Le point I est le milieu de [EF].

Dans toute la suite de l'exercice, on se place dans le repère orthonormé $\left(\text{A}~;~\vect{\text{AB}}~;~\vect{\text{AD}}~;~\vect{\text{AE}}\right)$.

Ainsi, par exemple, les points F{}, G et J ont pour coordonnées 

\begin{center}F(1~;~0~;~1),\qquad  G(1~;~1~;~1)\quad et\quad  J(2~;~0~;~1).\end{center}

\medskip

\begin{enumerate}
\item 




Puisque ABCDEFGH et BKLCFJMG sont des cubes on a FG = BF = EH = AD $ = 1$ et comme I est le milieu de [EF], on a FI $= \dfrac12 \times $ EF $ = \dfrac12 \times 1 = \dfrac12$.

En prenant comme base le triangle rectangle isocèle BFG et la hauteur [FI], on a donc :

$V_{\text{FIGB}} = \dfrac13 \times \dfrac12 \times 1 \times 1 \times \dfrac12 = \dfrac{1}{12}$~(unité de volume).
\item 
On a $\vect{\text{AI}} = \vect{\text{AE}} + \vect{\text{EI}} = \vect{\text{AE}} + \dfrac12\vect{\text{EF}} = \vect{\text{AE}} + \dfrac12\vect{\text{AB}} + 0\vect{\text{AD}}$. les coordonnées de I sont donc : $\left(\dfrac12~;~0~;~1\right)$.
\item 
Avec J(2~;~0~;~1) et D(0~;~1~;~0), on obtient $\vect{\text{DJ}}\begin{pmatrix}2\\-1\\1\end{pmatrix}$.

D'autre part $\vect{\text{BI}}\begin{pmatrix}- \dfrac12\\0 - 0\\1 - 0 \end{pmatrix}$ soit $\vect{\text{BI}}\begin{pmatrix}- \dfrac12\\0 \\1\end{pmatrix}$.

De G(1~;~1~;~1) on obtient $\vect{\text{BG}}\begin{pmatrix}0\\1\\1\end{pmatrix}$.

On a donc $\vect{\text{DJ}}\cdot \vect{\text{BI}} = -1 + 0 + 1 = 0$ et 

$\vect{\text{DJ}}\cdot \vect{\text{BG}} = 0 - 1 + 1 = 0$.

Conclusion $\vect{\text{DJ}}$ est un vecteur normal à deux vecteurs non colin"aires du plan (BIG) : c'est un vecteur normal à ce plan.
\item 
On sait donc d'après la question précédente que :

$M(x~;~y~;~z) \in (\text{BIG}) \iff 2x - 1y + 1z + d = 0$ avec $d \in \R$.

Or B(1~;~0~;~0)$ \in  (\text{BIG}) \iff 2\times 1 - 1 \times 0 + 1 \times 0 + d = 0 \iff 2 + d = 0 \iff d = - 2$.

On a donc $M(x~;~y~;~z) \in (\text{BIG}) \iff 2x - 1y + 1z - 2 = 0$
\item 
La droite $d$ orthogonale à (BIG) a un vecteur directeur colinéaire à $\vect{\text{DJ}}$.

On a F(1~;~0~;~1).D'autre part donc $M(x~;~y~;~z) \in d \iff$

$ \vect{\text{F}M}  = t \vect{\text{DJ}}, \: t \in \R \iff \left\{\begin{array}{l c l}
x - 1&=&2t\\
y - 0&=&- t\\
z- 1&=&t
\end{array}\right., \:t \in \R \iff \left\{\begin{array}{l c l}
x &=&1 + 2t\\
y - 0&=&\phantom{1 + }- t\\
z&=&1 + t
\end{array}\right., \:t \in \R$

\item
	\begin{enumerate}
		\item 
Si L$(x~;~y~;~z)$ est commun à $d$ et au plan (BIG) ses coordonnées vérifient les équations paramétriques de $d$ et l'équation du plan (BIG) soit le système :
		
$\left\{\begin{array}{l c l}
x &=&1 + 2t\\
y - 0&=&\phantom{1 + }- t\\
z&=&1 + t\\
2x - y + z - 2&=&0
\end{array}\right., \:t \in \R$.

En remplaçant $x, y$ et $z$ par leurs expressions en fonction de $t$ dans la dernière équation on obtient :

$2(1 + 2t) - (- t) + (1 + t) - 2 = 0 \iff 2 + 4t + t + 1 + t - 2 = 0\iff 6t + 1 = 0\iff t = - \dfrac16$ et en remplaçant dans $x, y$ et $z$, on obtient $x = 1 + \dfrac26 = \dfrac23$,\:$y = -\left(- \dfrac16 \right) = \dfrac16$ et $z = 1 - \dfrac16  = \dfrac56$.

Conclusion avec L $ = d \cap \text{(BIG)},\quad  \text{L}\left(\dfrac23~;~\dfrac16~;~\dfrac56\right)$.
		

		\item Calculer la longueur FL.
Avec $\vect{\text{FL}}\begin{pmatrix} -\frac13\\\frac16\\- \frac16 \end{pmatrix}$, 

on a FL$^2 = \left\|\vect{\text{FL}}\right\|^2 = \left(-\dfrac13\right)^2 + \left(\dfrac16\right)^2  + \left(- \dfrac16 \right)^2  = \dfrac19 + \dfrac{1}{36} + \dfrac{1}{36} =  \dfrac{4}{36} + \dfrac{2}{36} = \dfrac{6}{36}$,

donc FL $ = \sqrt{\dfrac{6}{36}} = \dfrac{\sqrt{6}}{6}$.
		\item 
En prenant comme base le triangle (BIG) le tétraèdre FIGB a pour hauteur [FL] ; on a donc 
		
$V_{\text{FIGB}} = \text{aire (BIG)} \times \text{[FL]} \times \dfrac13 \iff \dfrac{\dfrac{1}{12} \times 3}{\dfrac{\sqrt{6}}{6}} = \text{aire (BIG)} = \dfrac{\sqrt{6}}{4}$. (unités d'aire)
	\end{enumerate}
\end{enumerate}

\medskip
